\subsection{Prompts}
\label{app:prompts}

The rubric-instantiation prompt (\texttt{\{problem\}} is the task description):
\begin{quote}\small\begin{verbatim}
Here is a repository issue that a software-engineering agent will be asked to fix:

{problem}

Write an instance-specific rubric for judging an agent's trace on THIS issue:
for each section below, one sentence saying what specifically to look for in
this instance (name the likely relevant modules/files/tests if inferable).

Sections: understanding, localization, reproduction, editing, verification, final_state.
Respond with ONLY a JSON object with those keys, each a one-sentence string.
\end{verbatim}\end{quote}

The extraction prompt (\texttt{\{rubric\}} is the instantiated \trubric{} for the task, or the fixed generic rubric below for the generic baseline; \texttt{\{keys\}} the field names; the rendered trace is appended after the prompt):
\begin{quote}\small\begin{verbatim}
Below is the (truncated) execution trace of a software-engineering agent.
Extract a factual description of the agent's behavior for EACH of the following
sections. 30-60 words per section. Describe behavior only -- do not restate the
problem. If a section did not occur, write exactly what is missing (e.g. "never
ran any test after editing").

{rubric}

Respond with ONLY a JSON object with keys {keys}, each a string.
\end{verbatim}\end{quote}

The generic rubric:
\begin{quote}\small\begin{verbatim}
- understanding: how the agent oriented itself before acting
- localization: how it searched/navigated to find the relevant code
- reproduction: whether/how it reproduced the issue before fixing
- editing: what it changed (files, nature and size of the edit)
- verification: whether it re-ran tests or a repro script after editing, and the OBSERVED outcome (passing/failing/errors)
- final_state: how the run ended (clean finish, step/cost limit, submitted unverified)
\end{verbatim}\end{quote}
